%%%%%%%%%%% DISCLAIMER: THIS IS NOT AN OFFICIAL TUM TEMPLATE!!! %%%%%%%%%%%%%%%%%%%%%%
% Author: Jakob Thumm
% ----------------------- TODO ---------------------------
% Adapt these
\newcommand{\Chair}{Chair of something interesting but vague enough to get funding for everything}
\newcommand{\Student}{John Doe}
\newcommand{\Title}{Your title}
\newcommand{\Supervisor}{Prof. Dr.-Ing. Jane Doe}
\newcommand{\School}{School of Computation, Information and Technology}
% ----------------------- TODO ---------------------------

\documentclass[a4paper, 12pt]{scrartcl}

\usepackage[utf8]{inputenc}
\usepackage[ngerman]{babel}
\usepackage{amsmath}
\usepackage{amssymb}
\usepackage{fancyhdr}
\usepackage{color}
\usepackage{graphicx}
\usepackage{lastpage}
\usepackage{listings}
\usepackage{tikz}
\usepackage{pdflscape}
\usepackage{subfigure}
\usepackage{float}
\usepackage{polynom}
\usepackage{hyperref}
\usepackage{tabularx}
\usepackage{forloop}
\usepackage{geometry}
\usepackage{listings}
\usepackage{fancybox}
\usepackage{tikz}
\usepackage{algpseudocode,algorithm,algorithmicx}
\usepackage{cleveref}
\usepackage{lipsum}% just to generate text for the example

\let\origsection\section\renewcommand\section{\nostarsection}

\newcommand\nostarsection[1]{\sectionprelude\origsection{#1}\sectionpostlude}

\newcommand\sectionprelude{%
  \vspace{0.5em}
}

\newcommand\sectionpostlude{%
  \vspace{0em}
}

%Größe der Ränder setzen
\geometry{a4paper, left=12.7mm, right=15mm, top=33mm, bottom=12.7mm}
\usepackage{helvet}
\renewcommand{\familydefault}{\sfdefault}
\linespread{1.5}
\newlength\tindent
\setlength{\tindent}{\parindent}
\setlength{\parindent}{0pt}
\renewcommand{\indent}{\hspace*{\tindent}}
\setlength{\parskip}{0.6em}

%Kopf- und Fußzeile
\setlength{\headheight}{27.87993pt}
\pagestyle {fancy}
\fancyhf{} % sets both header and footer to nothing
\renewcommand{\headrulewidth}{0pt}
\definecolor{TUMalpha}{RGB}{151, 186, 218}
\fancyhead[L]{\textcolor{TUMalpha}{Exposé of \Student \\
\Chair}}
\fancyhead[R]{\includegraphics[height=0.8cm]{TUM_Logo_blau_cmyk_pdf.pdf}}

\fancyfoot[L]{}
\fancyfoot[C]{}
\fancyfoot[R]{Page \thepage /\pageref*{LastPage}}


\begin{document}
\,\\[1cm]
\begin{large}
Exposé \hfill \today\\[0.5cm]
Title of the disseratation (working title):\\[0.2cm]
\Title\\[2cm]
Doctoral candidate: \Student\\[0.5cm]
Supervisor: \Supervisor\\[0.5cm]
Academic department: \School\\[0.5cm]
Chair: \Chair\\[5cm]
\end{large}
\begin{tabularx}{\textwidth}{X X}
\rule{6cm}{0.5pt} & \rule{6cm}{0.5pt}\\
\textit{place, date} & \textit{place, date}\\[2cm]
\rule{6cm}{0.5pt} & \rule{6cm}{0.5pt}\\
Doctoral candidate & Supervisor
\end{tabularx}\\[1cm]
\textcolor{TUMalpha}{The formal examination takes place through the confirmation of the Graduate Center of the \School via DocGS. Please upload the exposé to DocGS.}
\newpage
\section{Introduction}
An exposé is a document in which you narrow down the scientific framework of your dissertation at an early stage and develop a realistic work plan. Since both the framework conditions and your own research can only be planned to a limited extent in scientific operations, the exposé will have to be modified regularly in the further course. Ideally, the exposé already provides the basis for the dissertation.
The length of the synopsis should be around 8-10 pages (12-point font; Arial; one and a half lines; plus images).

Structure of the exposé

\begin{enumerate}
    \item Introduction / Topic of the dissertation project
    \item Question
    \item Objective
    \item Research status (State of science and technology)
    \item Procedure and method (own attempts)
    \item Work and schedule
    \item Bibliography
    \item Appendix (if applicable)
\end{enumerate}
This structure is only intended as a suggestion and recommendation. 
In the letterhead from page 2 onward, please add Exposé of Max Mustermann und Chair of XXX.

\subsection{Reference}
Text Text Text Text Text Text Text Text Text Text Text Text

Text Text Text Text Text Text Text Text Text Text Text Text Text Text Text Text Text Text Text Text Text Text Text Text Text Text Text Text Text Text Text Text Text Text Text Text Text Text Text Text Text Text Text Text Text Text Text Text Text Text Text Text Text Text Text Text Text Text Text Text Text Text Text Text Text Text Text Text Text Text Text Text Text Text Text Text Text Text Text Text Text Text Text Text Text Text Text Text Text Text Text Text Text Text Text Text

\section{Question}
\lipsum[1]

\section{Objective}
\section{Research status}
\section{Procedure and method}
\section{Work and schedule}

%% Bibliography
\end{document}